\section{Conclusion}
\label{conclusion}

In this paper, we presented \appname, an agent coordination framework for automating defect impact analysis in DL libraries.
\appname tackles the key challenge of determining whether bug fixes in libraries like \textit{Transformers} and \textit{Megatron} affect downstream client programs.
By coordinating Miner, Diff Analyzer, and Impact Analyzer agents—\appname bridges the semantic gap between low-level library patches and high-level user-facing behavior, enabling context-aware impact assessment. 
Our evaluation on real-world PRs and commits demonstrates its effectiveness, achieving up to 95\% F1 in defect classification and 80\% precision in impact analysis.
In the future, we plan to extend its capabilities in several directions: (1) broadening support for additional DL libraries (e.g., vllm~\cite{vllm}) and defect types, (2) scaling the evaluation to larger datasets of client programs.
While the current workflow focuses on PRs and commits containing explicit fixes, future extensions may also consider report-only or reproduction-case scenarios.
% Notably, since CVE datasets for DL libraries are sparse, as most defects are fixed directly through PRs or commits, our evaluation reflects realistic workflows.
We further plan to extend \appname to security vulnerability analysis, enabling it can interpret CVE reports and perform structured impact reasoning similar to defect cases.